\documentclass[10pt,a4paper]{amsart}
\usepackage[margin=19mm]{geometry}
\usepackage{amsmath,amssymb,mathtools}
\usepackage[scr=rsfs,cal=euler]{mathalfa}
\usepackage{xfrac}
\usepackage[unicode,hidelinks,bookmarks=false]{hyperref}
\newcommand{\ie}{i.e.\ }
\newcommand{\cf}{cf.\ }
\newcommand{\resp}{respectively\ }
\newcommand{\Subsection}{\subsection}
\providecommand{\nobreakdash}{\nobreak\hbox{-}}
\setlength{\emergencystretch}{2em}
\multlinegap0pt
\usepackage{bm}
\usepackage{bbm}
\usepackage{dsfont}
\usepackage{xspace}
\usepackage{cancel}
\usepackage{mathtools}
\usepackage[shortlabels]{enumitem}
\usepackage{amsthm}
\makeatletter
\@ifundefined{Theorem}{\newtheorem{Theorem}{Theorem}}{}
\@ifundefined{Lemma}{\newtheorem{Lemma}[Theorem]{Lemma}}{}
\@ifundefined{Proposition}{\newtheorem{Proposition}[Theorem]{Proposition}}{}
\makeatother
\DeclareMathOperator{\supin}{in}
\def\strok{\restriction}
\newcommand{\Mon}{\mathfrak C}
\newcommand{\conv}{\mathrm{conv}}
\newcommand{\tp}{\mathrm{tp}}
\title[Countable models of weakly quasi-o-minimal theories II]{Countable models of weakly quasi-o-minimal theories II}
\author{Slavko Moconja, Predrag Tanovi\'c}
\date{Source: arXiv:2603.06430v1 (CC BY 4.0)}
\begin{document}
\maketitle

\noindent\emph{Source and licence.} Countable models of weakly quasi-o-minimal theories II. Version arXiv:2603.06430v1 (CC BY 4.0), distributed under the Creative Commons Attribution 4.0 International licence, as recorded in the arXiv licence metadata for this exact version and shown on the abstract page https://arxiv.org/abs/2603.06430v1. Full rights notice: licenses/arxiv-2603.06430-CC-BY-4.0.txt.\par\smallskip

\noindent\textit{Editorial scope.} Counted unit: the Proposition whose source label is \texttt{Prop\_quasio\_satisfy\_R} in the article (source0.tex lines 923--925, source label \texttt{Prop\_quasio\_satisfy\_R}), delivered together with the article's own complete original proof of it (source0.tex lines 926--958, one \texttt{proof} environment of 33 lines). No mathematics is written, repaired or re-derived: statement and proof are the article's own text line for line. Required context retained uncounted: Lemma O EKVIVALENCIJAMA (lines 462--469). Every definition, display and label of those lines is retained unchanged. Omitted from this package and not redistributed: 1--461, 470--922, 959--1610. Nothing inside a retained excerpt is omitted or altered: every retained body is delivered exactly as the article's lines are written, so no omission receipt of any kind accompanies this package. Citations: the retained excerpts carry no citation command at all, so no bibliography entry of the article is transcribed and no reference is redistributed. Reference adaptation: none was needed and none was made; every reference inside the delivered unit resolves inside it, so no unresolved reference or citation is produced. Printed slips kept literal: no printed word, symbol, hypothesis or number of the counted statement or of its proof was altered. Layout and class: the article's own class is \texttt{amsart} with packages including amssymb, amsmath, amsthm, enumitem, eucal, mathabx, xcolor, hyperref, biblatex, etoolbox, orcidlink, geometry; the wrapper is the lane's pinned \texttt{amsart} editorial wrapper at 10pt on A4, which restates the theorem-like environments the retained text needs and adds the article's own macro definitions that the retained text needs (\texttt{\textbackslash Mon}, \texttt{\textbackslash conv}, \texttt{\textbackslash strok}, \texttt{\textbackslash supin}, \texttt{\textbackslash tp}), with extra wrapper packages (bm, bbm, dsfont, xspace, cancel, mathtools, enumitem). The delivered numbering is the unit's own. Assets: no retained span contains a figure environment, a graphics-inclusion command, a PDF-page inclusion, a table or a TikZ picture; no image file is copied into this package, the package declares no asset dependency and no image file is redistributed with it. Licence: version arXiv:2603.06430v1 of this article is distributed under the Creative Commons Attribution 4.0 International licence, as recorded in the arXiv licence metadata for this exact version and shown on the abstract page https://arxiv.org/abs/2603.06430v1. A full rights notice is delivered at licenses/arxiv-2603.06430-CC-BY-4.0.txt. Characters: every retained excerpt is pure ASCII, so the delivered engine, XeLaTeX with its default Latin Modern text font, reports no missing character.
\begin{Lemma}\label{Lemma O EKVIVALENCIJAMA}
Let $A\subseteq B$, let $p\in S(A)$ be a weakly o-minimal type, and let $q\in S(B)$ be its extension. Suppose that $E\in\mathcal E_q\smallsetminus\{\mathbf 1_q\}$. The following conditions are equivalent:
\begin{enumerate}[\hspace{10pt}1)]
\item There exists $E_0\in\mathcal E_p$ such that $(E_{0})_{\strok q(\Mon)}= E$;
\item $E$ is relatively $A$-definable;
\item For some (all) $a\models q$, the class $[a]_E$ is $Aa$-invariant.
\end{enumerate}
\end{Lemma}

\begin{Proposition}\label{Prop_quasio_satisfy_R} 
Quasi-o-minimal theories satisfy condition (R).
\end{Proposition}

\begin{proof}
Suppose that $T$ is quasi-o-minimal with respect to $<$, and let $p_A\in S_1(A)$ and $E_A\in\mathcal E_p$, and we prove that $E_A$ is relatively 0-definable. We may assume $E_A\neq\mathrm{id}_{p_A(\Mon)}$ and $E_A\neq \mathbf 1_p$. 
Let $p_0=(p_A)_{\strok\emptyset}$, and fix $a\models p_A$.
Since $p_A(\Mon)$ is a convex subset of $p_0(\Mon)$ and $[a]_E$ is a bounded and convex subset of $p_A(\Mon)$, $[a]_E$ is a bounded and convex subset of $p_0(\Mon)$, and it is relatively $Aa$-definable within $p_0(\Mon)$. We prove that there are $\theta(x)\in p_0(x)$ and $b,c\in\Mon$ such that:
\setcounter{equation}{0}
\begin{equation}
[a]_{E_A}=(\theta(\Mon)\cap (b,c))\cap p_0(\Mon).
\end{equation}
As $T$ is quasi-o-minimal, the formula relatively defining $[a]_{E_A}$ within $p_0(\Mon)$ can be chosen in the form $\bigvee_{i=1}^n\theta_i(x)\land x\in J_i$, where $\theta_i(x)$ are $L$-formulae and $J_i$ are intervals in $\Mon$. Since $[a]_{E_A}\subset p_0(\Mon)$, we may assume that $\theta_i(x)\in p_0(x)$ for all $i=1,\dots,n$, and note that $\theta(x)\land\bigvee_{i=1}^nx\in J_i$ still relatively defines $[a]_{E_A}$ within $p_0(\Mon)$, where $\theta(x)=\bigwedge_{i=1}^n\theta_i(x)$. Also, we may assume that each $J_i$ meets $[a]_{E_A}$. The left end, say $b$, of some $J_i$ satisfies $b\leqslant[a]_{E_A}$, and the right end, say $c$, of some $J_i$ satisfies $[a]_{E_A}\leqslant c$. It is easy to see that $[a]_{E_A}$ is relatively defined by $\theta(x)\land x\in J$ within $p_0(\Mon)$, where $J$ is some interval with endpoints $b$ and $c$. 
Now, note that the class $[a]_{E_A}$ does not have a minimal (maximal) element; otherwise, every realization of $p_A$ would be minimal (maximal) in its $E_A$-class, so we would have $E_A=\mathrm{id}_{p_A(\Mon)}$, which is not the case. Thus $b,c\notin[a]_{E_A}$, so we may assume that $J=(b,c)$, and we have established (1). 

By Lemma \ref{Lemma O EKVIVALENCIJAMA} it suffices to prove that the class $[a]_{E_A}$ is $a$-invariant. Toward a contradiction, suppose that there is $A'\equiv A\,(a)$ such that $[a]_{E_A}\neq [a]_{E_{A'}}$ (where $E_{A'}$ is the conjugate of $E_A$). Then, after possibly reversing the order and/or exchanging the roles of $A$ and $A'$, we can assume $\supin_{p_0(\Mon)} ([a]_{E_{A'}})\subset\supin_{p_0(\Mon)} ([a]_{E_A})$. 
%
%Let $\phi(x)$ be a relative definition of the class $[a]_{E_A}$ within $p_0(\Mon)$, where $p_0=(p_A)_{\strok \emptyset}$. We prove that $\phi(x)$ can be chosen such that $\phi(\Mon)=D\cap (b,c)$, where $D$ is 0-definable and $b,c\in\Mon$. By quasi-o-minimality, we can assume that $\phi(x)$ is a Boolean combination of $L$-formulas $\theta_1(x),\dots,\theta_n(x)$ and formulas of the form $x<d$ or $x\leqslant d$ ($d\in\Mon$); without loss assume that $\phi(x)$ is in the disjunctive normal form. Observe that for each $i\leqslant n$ exactly one of $\theta_i(x)$ and $\lnot \theta_i(x)$ is consistent with $p_0(x)$, so there is a unique vector $(\epsilon_1,\dots,\epsilon_n)\in 2^{[n]}$ such that $\bigwedge_{i\in[n]}\theta_i^{\epsilon_i}(x)$ is consistent with $p_0(x)$; denote this conjunction by $\theta(x)$. Then $\phi(x)$ is a Boolean combination of $\theta(x)$ and formulas of the form $x<d$ or $x\leqslant d$ ($d\in\Mon$). 
%Since $[a]_{E_A}$ is a convex subset of $p_0(\Mon)$,  
%$\phi(x)$ can be chosen to define the set $D\cap I$, where $D=\theta(\Mon)$ and $I$ is an interval with endpoints $b$ and $c$ ($b<c$). Note that $b=-\infty$ or $c=+\infty$ is impossible, since $[a]_{E_A}$ is a bounded subset of $p_0(\Mon)$ and $p_0(\Mon)\subseteq D$. Thus, $b,c\in\Mon$. 
%Also note that $b\notin [a]_{E_A}$ and $c\notin [a]_{E_A}$ because $[a]_{E_A}$ does not have a minimal (maximal) element. Therefore:\setcounter{equation}{0}
%\begin{equation}
%[a]_{E_A}=(D\cap (b,c))\cap p_0(\Mon).
%\end{equation}
%
%
Denote: \ $q_0=\tp(c)$, $q_A=\tp(c/A)$,  and $C_A(a)=\{x\in p_0(\Mon)\mid [a]_{E_A}<x\}$.  From (1) we derive:
\begin{equation}
[a]_{E_A}<c \ \text{ and } \  c\leqslant C_A(a).
\end{equation}
Note that the set $C_A(a)\cap p_A(\Mon)$ is an initial part of $C_A(a)$ that is closed for $E_A$-classes. Consequently, since no $E_A$-class contains a minimal element, neither does $C_A(a)$, so $c\notin C_A(a)$. Then $[a]_{E_A}<c<C_A(a)$ implies $c\notin p_0(\Mon)$, so $q_0\neq p_0$.

Choose $a_0c_0\equiv ac\,(A)$ with $[a_0]_{E_A}<[a]_{E_A}$. Then $c_0\models q_A$, $[a_0]_{E_A}<c_0$, and there are no realizations of $p_A$ between $[a_0]_{E_A}$ and $c_0$. Since $q_A\neq p_A$ and since $[a_0]_{E_A}<[a]_{E_A}$, we conclude $c_0<[a]_{E_A}<c$. Next, choose $c'\models q_0$ so that $c'A'\equiv cA\,(a)$. 
Note that $\supin_{p_0(\Mon)} ([a]_{E_{A'}})\subset\supin_{p_0(\Mon)} ([a]_{E_A})$ implies that $[a]_{E_{A'}}<d$ for some $d\in [a]_{E_A}$. 
Then $d\in C_{A'}(a)$, so $c'<d$ and $c'\in \conv([a]_{E_A})$ (the convex hull of $[a]_{E_A}$ in $(\Mon,<)$). 
We conclude that $c'$ divides the set $\conv([a]_{E_A})$ into two parts, each containing realizations of $p_A(\Mon)$. Also, $c'\in \conv([a]_{E_A})$ combined with $c_0<[a]_{E_A}<c$ yields $c_0<c'<c$. Since $q_A(\Mon)$ is a convex subset of $q_0(\Mon)$, $c_0<c'<c$ together with $c_0,c\models q_A$ and $c'\models q_0$ implies $c'\models q_A$; hence $c\equiv c'\,(A)$. By our choice of $c$, $[a]_{E_A}<c<C_A(a)$ implies that $c$ does not divide the convex hull of any $E_A$-class; on the other hand, $c'$ divides $\conv([a]_{E_A})$. This contradicts $c\equiv c'\,(A)$ and completes the proof of the proposition. 
\end{proof}

\end{document}
